% Generated draft: requires mathematical and source review.
\documentclass[11pt]{article}
\usepackage[margin=25mm]{geometry}
\usepackage{fontspec}
\setmainfont{Latin Modern Roman}
\usepackage{amsmath,amssymb,mathtools,unicode-math}
\setmathfont{Latin Modern Math}
\usepackage{enumitem,needspace,xurl,hyperref}
\setlength{\parindent}{0pt}
\setlength{\parskip}{5pt}
\setlength{\emergencystretch}{3em}
\newcommand{\N}{\mathbb{N}}
\newcommand{\Z}{\mathbb{Z}}
\newcommand{\Q}{\mathbb{Q}}
\newcommand{\R}{\mathbb{R}}
\newcommand{\C}{\mathbb{C}}
\newcommand{\F}{\mathbb{F}}
\newcommand{\A}{\mathbb{A}}
\newcommand{\PP}{\mathbb{P}}
\newcommand{\E}{\mathbb{E}}
\newcommand{\eps}{\varepsilon}
\newcommand{\dd}{\,\mathrm{d}}
\begin{document}
\section*{UnsolvedMath ID: 20003173}
\textbf{Source problem code:} catalogue entry
\subsection{Definitions and mathematical statement}
Let $(M^n,g(t))$, $n\geq4$, be a Ricci flow from a compact boundaryless Riemannian manifold with positive isotropic curvature (PIC). With a consistent curvature-tensor convention, PIC means
\[
 R_{1313}+R_{1414}+R_{2323}+R_{2424}-2R_{1234}>0
\]
for every orthonormal four-frame. Nonnegative curvature operator (NCO) means $\langle\operatorname{Rm}\omega,\omega\rangle\geq0$ for every $\omega\in\Lambda^2TM$. The source asks whether every blow-up limit has NCO.

The supplied assessment concerns standard finite-time singularity-model blow-ups: for $t_i\to T<\infty$ and curvature scales $Q_i\to\infty$, take pointed smooth limits of
\[
 g_i(s)=Q_i g(t_i+s/Q_i),
\]
with completeness, bounded curvature on the required time intervals, and the ancient/noncollapsing hypotheses used by the model-classification results. It reports an affirmative answer in dimension $4$ and every $n\geq12$ from published pinching and ancient-$\kappa$-solution classification results. Assuming Zhengnan Chen's 2025 preprint theorem that complete ancient uniformly PIC flows in dimensions $n\geq9$ are weakly PIC2, it reports the same conclusion conditionally for $n=9,10,11$. The checked theory in the supplied aid leaves $n=5,6,7,8$ open. Full references and exact model hypotheses are not supplied.

Here weak PIC2 means that the curvature tensor extended by two flat directions has nonnegative isotropic curvature. A uniform PIC condition is a positive, scale-invariant quantitative margin, such as membership of $\operatorname{Rm}-\theta\operatorname{scal}I$ in the nonnegative PIC cone for one fixed $\theta>0$, in the theorem's convention. These conditions must not be identified with NCO: the aid explicitly requires a separate compact/noncompact ancient-model classification bridge from the pinched cone to NCO.

The attempt's independent general lemma is a cone-defect limit statement. Let $\mathcal C$ be a closed frame-invariant curvature cone invariant under multiplication by positive scalars. Suppose, on the region being rescaled, scalar curvature $r\geq0$ and
\[
 \operatorname{Rm}+f(r)I\in\mathcal C,
\]
where $I$ is the constant-curvature identity operator and $f:[0,\infty)\to[0,\infty)$ is locally bounded with $f(r)/r\to0$ as $r\to\infty$. If a smooth parabolic blow-up has locally bounded rescaled scalar curvature, every limit curvature tensor lies in $\mathcal C$. Indeed for each $\eps>0$, $f(r)\leq\eps r+C_\eps$, so on a rescaled compact region with scalar curvature bounded by $S$ the error is
\[
 Q_i^{-1}f(r_i)\leq\eps S+C_\eps/Q_i\longrightarrow0
 \quad\text{after letting }\eps\to0.
\]
Closedness then removes the defect. Local boundedness is needed to control bounded original scalar values. This lemma yields NCO directly only when $\mathcal C$ is the NCO cone; for weak PIC2 it gives that weaker condition and still needs the stated classification bridge. It does not by itself settle the unresolved dimensions or every possible interpretation of blow-up limit.
\subsection{Short English statement}
Must singularity blow-ups of a compact PIC Ricci flow have nonnegative curvature operator? The supplied assessment answers yes for standard models in dimension $4$ and at least $12$, gives a preprint-dependent route for $9$--$11$, and leaves $5$--$8$ unresolved. A cone-defect lemma explains why sublinear curvature errors disappear under smooth rescaling, while preserving the distinction between weak PIC2 and NCO.
\subsection{Sources}
\begin{itemize}
\item Local UnsolvedMath registry, permanent ID $20003173$, source code catalogue entry; canonical location \texttt{aim-topology-notes.json}, \texttt{notes[260]}.
\item American Institute of Mathematics, \emph{Manifolds with nonnegative sectional curvature}, Problem 48, \url{https://aimath.org/WWN/nnsectcurvature/nnsectcurvature.pdf}. The supplied verification restores the lost superscript in $M^n$ and reports no other correction.
\item Zhengnan Chen, 2025 preprint theorem on ancient uniformly PIC flows, as described in the research summary; no title or identifier is supplied. Other published pinching and ancient-$\kappa$ classification inputs are mentioned without bibliographic details.
\item Research attempt 1 is marked partial result, proved, and low novelty confidence, with an explicit machine-generated, non-peer-reviewed research-aid caveat requiring independent verification before citation. The dimension-$9$--$11$ conclusion is expressly preprint-dependent.
\end{itemize}
\end{document}
